\chapter[Izomorfizmy i elementarność struktur]{Homomorfizmy, izomorfizmy i elementarność struktur}
\label{ch:model-morphisms}

\section{Dlaczego porównujemy struktury, a nie tylko ich uniwersa?}
W rozdziale~\ref{ch:semantics-fol} zdefiniowaliśmy strukturę jako zbiór
razem z interpretacjami symboli stałych, funkcyjnych i relacyjnych.
Dwa zbiory mogą mieć tyle samo elementów, a jednak pełnić zupełnie inne
role jako struktury. Na przykład na tym samym zbiorze $\{0,1\}$
można interpretować jednoargumentowy symbol relacyjny $R$
jako relację pustą albo jako relację pełną. Będą to różne struktury, mimo że ich uniwersa są identyczne.

Chcemy rozróżnić cztery pytania. Czy funkcja między uniwersami
zachowuje podstawowe działania? Czy zachowuje także relacje?
Czy pozwala odtworzyć strukturę bez utraty informacji?
Czy nawet wszystkie formuły pierwszego rzędu mają takie samo znaczenie?
Odpowiedzi prowadzą do homomorfizmów, izomorfizmów i odwzorowań elementarnych.

\section{Homomorfizm i zanurzenie struktury}
\begin{uwaga}[Jak czytamy $h:\mathfrak A\to\mathfrak B$?]
Zapis $h:\mathfrak A\to\mathfrak B$ jest skrótem oznaczającym
funkcję $h:A\to B$ między \emph{uniwersami} struktur, która zachowuje
wymagane interpretacje symboli. Nie oznacza, że struktury są argumentami
funkcji $h$. To skrócenie będzie używane w dalszych twierdzeniach.
\end{uwaga}

\begin{definicja}[Homomorfizm struktur]\label{def:model-homomorphism}
Niech $\mathcal L$ będzie ustalonym językiem pierwszego rzędu,
a $\mathfrak A,\mathfrak B$ strukturami dla $\mathcal L$
o uniwersach odpowiednio $A,B$. Funkcję $h:A\to B$ nazywamy
\emph{homomorfizmem struktur}, jeśli spełnia następujące warunki:
\begin{enumerate}[label=(\roman*)]
\item dla każdego symbolu stałej $c$ zachodzi $h(c^{\mathfrak A})=c^{\mathfrak B}$;
\item dla każdego $n$-argumentowego symbolu funkcyjnego $f$ i wszystkich
$a_1,\ldots,a_n\in A$ zachodzi
\[
h(f^{\mathfrak A}(a_1,\ldots,a_n))
=f^{\mathfrak B}(h(a_1),\ldots,h(a_n));
\]
\item dla każdego $n$-argumentowego symbolu relacyjnego $R$
i wszystkich $a_1,\ldots,a_n\in A$ zachodzi implikacja
\[
(a_1,\ldots,a_n)\in R^{\mathfrak A}
\quad\Longrightarrow\quad
(h(a_1),\ldots,h(a_n))\in R^{\mathfrak B}.
\]
\end{enumerate}
Homomorfizm nazywamy \emph{zanurzeniem}, jeśli ponadto jest iniekcją
i w ostatnim warunku zachodzi równoważność zamiast samej implikacji.
\end{definicja}

\begin{uwaga}[Dlaczego relacje traktujemy inaczej niż funkcje?]
Warunek dla funkcji jest równością wartości: obie kolejności
„najpierw działanie, potem $h$” oraz „najpierw $h$, potem działanie”
prowadzą do tego samego wyniku. Dla relacji homomorfizm zapewnia
jedynie zachowanie faktu, że pewna relacja zachodzi. Sam homomorfizm
nie musi zachowywać niezachodzenia relacji. Dlatego zanurzenie
wymaga także implikacji w drugą stronę.
\end{uwaga}

\begin{przyklad}[Zachowanie dodawania]
Rozważmy język z symbolem $+$ arności dwa i stałą $0$.
Funkcja $h:\Z\to\Z$, $h(n)=2n$, jest homomorfizmem
struktur $(\Z,+,0)$ do $(\Z,+,0)$, ponieważ
$h(m+n)=2(m+n)=2m+2n=h(m)+h(n)$ i $h(0)=0$.
Jest ona iniekcją, więc w tym języku, pozbawionym
symboli relacyjnych poza równością, jest zanurzeniem.
Nie jest izomorfizmem na całe $\Z$, gdyż liczba $1$
nie jest obrazem żadnego argumentu.
\end{przyklad}

\begin{przyklad}[Homomorfizm może gubić informację]
W języku z jednym dwuargumentowym symbolem relacyjnym $R$ niech
$A=\{a,b\}$, $B=\{c\}$ i niech w obu strukturach
$R$ oznacza relację pełną. Stała funkcja $h(a)=h(b)=c$
jest homomorfizmem: każdy fakt $R(x,y)$ prawdziwy
w pierwszej strukturze pozostaje prawdziwy po zastosowaniu $h$.
Nie jest jednak zanurzeniem, ponieważ nie jest iniekcją:
dwa różne obiekty zostały utożsamione.
\end{przyklad}

\begin{lemat}[Zachowanie wartości termów]\label{lem:model-homo-terms}
Niech $\mathcal L$ będzie językiem pierwszego rzędu,
$h:\mathfrak A\to\mathfrak B$ homomorfizmem struktur języka
$\mathcal L$, $s$ wartościowaniem w $\mathfrak A$,
a $t$ termem języka $\mathcal L$.
Niech $h\circ s$ oznacza wartościowanie $x\mapsto h(s(x))$ w $\mathfrak B$.
Wtedy
\[
h\bigl(t^{\mathfrak A}[s]\bigr)=t^{\mathfrak B}[h\circ s].
\]
\end{lemat}
\begin{proof}
Stosujemy indukcję po budowie termu. Jeśli $t=x$ jest zmienną,
lewa strona wynosi $h(s(x))$, co jest definicją prawej strony.
Jeśli $t=c$ jest stałą, równość wynika z pierwszego warunku
homomorfizmu. Niech teraz $t=f(t_1,\ldots,t_n)$.
Wartość lewej strony można zapisać jako
\[
h\bigl(f^{\mathfrak A}(t_1^{\mathfrak A}[s],\ldots,
 t_n^{\mathfrak A}[s])\bigr).
\]
Drugi warunek homomorfizmu pozwala przenieść $h$ do każdego argumentu.
Z założenia indukcyjnego argumenty te są równe
$t_i^{\mathfrak B}[h\circ s]$, więc otrzymujemy
$f^{\mathfrak B}(t_1^{\mathfrak B}[h\circ s],\ldots,t_n^{\mathfrak B}[h\circ s])$.
To dokładnie wartość termu $t$ po prawej stronie tezy.
\end{proof}


\section{Dokładnie co zachowuje homomorfizm?}
\begin{definicja}[Formuła dodatnia egzystencjalna]\label{def:model-positive-existential}
Niech $\mathcal L$ będzie językiem pierwszego rzędu.
Formułę języka $\mathcal L$ nazywamy \emph{dodatnią egzystencjalną},
jeśli otrzymano ją z formuł atomowych tego języka przez skończone
stosowanie wyłącznie koniunkcji $\land$, alternatywy $\lor$
i kwantyfikatora egzystencjalnego $\exists$.
Nie używamy negacji, implikacji ani kwantyfikatora ogólnego.
Na przykład $\exists y\,(R(x,y)\land S(y))$ jest taką formułą,
ale $\neg R(x,x)$ już nie.
\end{definicja}
\begin{twierdzenie}[Zachowanie formuł dodatnich egzystencjalnych]\label{thm:model-homo-positive}
Niech $h:\mathfrak A\to\mathfrak B$ będzie homomorfizmem struktur
tego samego języka. Jeśli $\varphi$ jest formułą dodatnią
egzystencjalną i $s$ wartościowaniem w $\mathfrak A$, to
\[
\mathfrak A\models\varphi[s]\quad\Longrightarrow\quad
\mathfrak B\models\varphi[h\circ s].
\]
Nie twierdzimy, że implikacja odwrotna zachodzi.
\end{twierdzenie}
\begin{proof}
Dla atomu $t=u$ używamy Lematu~\ref{lem:model-homo-terms}:
z równości wartości w $A$ wynika równość ich obrazów w $B$.
Dla atomu $R(t_1,\ldots,t_n)$ otrzymujemy obrazy wartości termów,
a warunek (iii) w definicji homomorfizmu zapewnia zachowanie relacji.
Jeżeli formuła ma postać $\psi\land\theta$, to obie formuły składowe
zachodzą w $\mathfrak A$; z założenia indukcyjnego obie zachodzą
w $\mathfrak B$. Przy $\psi\lor\theta$ wystarczy zastosować
założenie indukcyjne do jednego spełnionego składnika.
Wreszcie niech $\varphi=\exists x\,\psi$ i niech $a\in A$
będzie świadkiem. Otrzymujemy
$\mathfrak A\models\psi[s[x\mapsto a]]$.
Założenie indukcyjne daje
$\mathfrak B\models\psi[(h\circ s)[x\mapsto h(a)]]$;
a więc $h(a)\in B$ jest świadkiem dla formuły egzystencjalnej
w $\mathfrak B$. To wyczerpuje wszystkie dopuszczalne konstrukcje.
\end{proof}
\begin{przyklad}[Dlaczego nie wolno dopisać negacji]
W języku z jednoargumentową relacją $R$ weźmy
$A=\{a\}$, $B=\{b\}$, $R^{\mathfrak A}=\varnothing$
oraz $R^{\mathfrak B}=\{b\}$. Funkcja $h(a)=b$ jest
homomorfizmem: w $A$ nie zachodzi żaden atom $R(a)$,
więc nie ma faktu relacyjnego, który trzeba by zachować.
Wybierzmy zmienną $x$ i wartościowanie $s$ w $\mathfrak A$
spełniające $s(x)=a$. Wówczas
$\mathfrak A,s\models\neg R(x)$, ponieważ $a\notin R^{\mathfrak A}$,
ale $\mathfrak B,h\circ s\not\models\neg R(x)$,
ponieważ $h(a)=b\in R^{\mathfrak B}$.
Homomorfizm nie musi więc zachowywać negacji.
\end{przyklad}
\begin{przyklad}[Dlaczego nie wolno dopisać $\forall$]
Niech język zawiera jedną jednoargumentową relację $R$.
Weźmy $A=\{a\}$, $B=\{a,b\}$ oraz
$R^{\mathfrak A}=\{a\}$, $R^{\mathfrak B}=\{a\}$.
Włączenie $A\hookrightarrow B$ jest homomorfizmem,
a nawet zanurzeniem. W mniejszej strukturze zachodzi
$\forall x\,R(x)$, a w większej nie, gdyż $b\notin R^{\mathfrak B}$.
Wniosek: sama iniekcyjność homomorfizmu nie zastępuje
surjektywności ani elementarności.
\end{przyklad}

\section{Izomorfizmy i przenoszenie prawdziwości}
\begin{definicja}[Izomorfizm]\label{def:model-isomorphism}
Niech $\mathfrak A$ i $\mathfrak B$ będą strukturami tego samego
języka o uniwersach odpowiednio $A$ i $B$.
Izomorfizmem tych struktur nazywamy bijektywne zanurzenie $h:A\to B$. Piszemy wtedy
$\mathfrak A\cong\mathfrak B$. Struktury izomorficzne różnią się
jedynie nazwami elementów, nie zaś wynikami działań i relacjami.
\end{definicja}

\begin{przyklad}[Dwie nazwy tych samych obiektów]
Niech w języku z jedną dwuargumentową relacją $R$
struktura $\mathfrak A$ ma uniwersum $\{1,2\}$
i relację $R^{\mathfrak A}=\{(1,2)\}$, natomiast
$\mathfrak B$ ma uniwersum $\{a,b\}$ i
$R^{\mathfrak B}=\{(a,b)\}$. Funkcja $h(1)=a$,
$h(2)=b$ jest izomorfizmem. Przenosi jedyną parę
pozostającą w relacji na jedyną taką parę w drugiej strukturze.
Nie zmienia się żadna własność, którą można wyrazić formułą
rozważanego języka.
\end{przyklad}

\begin{twierdzenie}[Izomorfizm zachowuje wszystkie formuły]\label{thm:model-iso-invariance}
Niech $h:\mathfrak A\to\mathfrak B$ będzie izomorfizmem
struktur języka $\mathcal L$, a $\varphi$ dowolną formułą
tego języka. Dla każdego wartościowania $s$ w $\mathfrak A$ zachodzi
\[
\mathfrak A\models\varphi[s]
\quad\Longleftrightarrow\quad
\mathfrak B\models\varphi[h\circ s].
\]
W szczególności dla każdego zdania $\sigma$ zachodzi
$\mathfrak A\models\sigma\iff\mathfrak B\models\sigma$.
\end{twierdzenie}
\begin{proof}
Przeprowadzimy indukcję po budowie formuły.
Dla atomu $t=u$ z Lematu~\ref{lem:model-homo-terms} otrzymujemy
$h(t^{\mathfrak A}[s])=t^{\mathfrak B}[h\circ s]$
i analogiczną równość dla $u$. Ponieważ $h$ jest iniekcją,
dwie wartości w $A$ są równe wtedy i tylko wtedy, gdy ich obrazy
przez $h$ są równe. To daje tezę dla atomu z równością.
Dla atomu relacyjnego $R(t_1,\ldots,t_n)$ korzystamy
z Lematu~\ref{lem:model-homo-terms} oraz z faktu,
że zanurzenie zarówno zachowuje, jak i odzwierciedla relacje.

Negacja działa w obu strukturach tak samo: $\neg\psi$ jest
spełniona dokładnie wtedy, gdy $\psi$ nie jest spełniona.
Stosując założenie indukcyjne do $\psi$, otrzymujemy tezę
dla $\neg\psi$. Dla każdego spójnika dwuargumentowego
($\land$, $\lor$, $\to$, $\leftrightarrow$) jego tabela prawdy
zależy wyłącznie od wartości logicznych składników,
więc założenie indukcyjne daje identyczny wynik po obu stronach.

Niech $\varphi=\exists x\,\psi$. Jeśli
$\mathfrak A\models(\exists x\,\psi)[s]$,
istnieje $a\in A$ z $\mathfrak A\models\psi[s[x\mapsto a]]$.
Na mocy założenia indukcyjnego otrzymujemy
$\mathfrak B\models\psi[(h\circ s)[x\mapsto h(a)]]$,
a więc $\mathfrak B\models(\exists x\,\psi)[h\circ s]$.
Odwrotnie, jeśli istnienie jest spełnione w $\mathfrak B$,
możemy wybrać tam świadka $b\in B$. Surjekcja $h$ daje
$a\in A$ z $h(a)=b$, więc przez założenie indukcyjne
otrzymujemy świadka również w $\mathfrak A$.
Kwantyfikator ogólny można rozpatrzyć analogicznie,
korzystając z tego, że bijekcja $h$ przebiega wszystkie elementy
$B$ dokładnie raz; albo użyć równoważności
$\forall x\,\psi\iff\neg\exists x\,\neg\psi$.
Wszystkie formy budowy formuły zostały sprawdzone.
\end{proof}

\section{Podstruktury i to, czego nie zachowują}\label{sec:model-substructure}
\begin{definicja}[Podstruktura]\label{def:model-substructure}
Niech $\mathcal L$ będzie językiem pierwszego rzędu,
a $\mathfrak A,\mathfrak B$ strukturami tego języka
o niepustych uniwersach $A\subseteq B$.
Mówimy, że $\mathfrak A$ jest
\emph{podstrukturą} $\mathfrak B$ i piszemy
$\mathfrak A\subseteq\mathfrak B$, gdy dla każdego symbolu
stałej $c$ zachodzi $c^{\mathfrak A}=c^{\mathfrak B}\in A$,
dla każdego $n$-argumentowego symbolu funkcyjnego $f$
zbiór $A$ jest zamknięty na $f^{\mathfrak B}$ oraz
$f^{\mathfrak A}=f^{\mathfrak B}|_{A^n}$,
a dla każdego $n$-argumentowego symbolu relacyjnego $R$
zachodzi $R^{\mathfrak A}=R^{\mathfrak B}\cap A^n$.
\end{definicja}

\begin{przyklad}[Podstruktura, która nie zachowuje zdania]
W języku z jednoargumentowym symbolem funkcyjnym $S$
rozważmy $\mathfrak A=(\N,S)$ i $\mathfrak B=(\Z,S)$,
gdzie $S(n)=n+1$ w obu strukturach. Zbiór $\N$
jest zamknięty na działanie $S$, więc
$\mathfrak A\subseteq\mathfrak B$.
W $\mathfrak B$ prawdziwe jest zdanie
$\forall x\,\exists y\,(S(y)=x)$: każdy $x\in\Z$
ma poprzednik $y=x-1$. W $\mathfrak A$ zdanie jest fałszywe,
gdyż element $0$ nie ma poprzednika w $\N$.
Bycie podstrukturą nie gwarantuje więc zachowania
wszystkich zdań z kwantyfikatorem egzystencjalnym.
\end{przyklad}

\begin{definicja}[Zdanie uniwersalne o macierzy bez kwantyfikatorów]
Niech $\mathcal L$ będzie językiem pierwszego rzędu.
Zdaniem \emph{uniwersalnym o macierzy bez kwantyfikatorów}
nazywamy zdanie języka $\mathcal L$ postaci
$\forall x_1\cdots\forall x_n\,\psi(x_1,\ldots,x_n)$,
gdzie $n\geq 0$, zmienne $x_1,\ldots,x_n$ są parami różne,
wszystkie zmienne wolne formuły $\psi$ należą do tej listy,
a $\psi$ nie zawiera kwantyfikatorów. Gdy $n=0$,
formuła $\psi$ jest zdaniem bez kwantyfikatorów.
\end{definicja}

\begin{twierdzenie}[Zachowanie zdań uniwersalnych]\label{thm:model-universal-substructure}
Niech $\mathfrak A\subseteq\mathfrak B$ będą strukturami
tego samego języka $\mathcal L$, a $\sigma$ zdaniem
uniwersalnym o macierzy bez kwantyfikatorów tego języka.
Jeśli $\mathfrak B\models\sigma$, to $\mathfrak A\models\sigma$.
\end{twierdzenie}
\begin{proof}
Niech $\sigma=\forall x_1\cdots\forall x_n\,\psi$.
Wybierzmy dowolne $a_1,\ldots,a_n\in A$.
Ponieważ $A\subseteq B$, można w $\mathfrak B$
przypisać zmiennym te same wartości. Z założenia
$\mathfrak B\models\sigma$ wynika, że
$\psi$ jest spełniona w $\mathfrak B$ dla tych wartości.
Dla termów interpretacje w $\mathfrak A$ i $\mathfrak B$
są identyczne, gdy wszystkie argumenty leżą w $A$:
dowodzi się tego indukcją po budowie termu, używając
definicji podstruktury. Dlatego formuły atomowe
zachowują identyczną wartość logiczną w obu strukturach;
jeśli $R$ jest symbolem relacyjnym arności $k$,
korzystamy z równości
$R^{\mathfrak A}=R^{\mathfrak B}\cap A^k$.
Indukcja po budowie kwantyfikatorowo wolnej $\psi$
przenosi tę zgodność przez negację i wszystkie spójniki.
Zatem $\psi$ jest spełniona w $\mathfrak A$
przy dowolnych $a_1,\ldots,a_n\in A$,
czyli $\mathfrak A\models\sigma$.
Jeśli $n=0$, rozumowanie sprowadza się do zgodności
prawdziwości zamkniętej formuły bez kwantyfikatorów;
nie trzeba wtedy wybierać żadnych parametrów.
\end{proof}

\section{Podstruktura elementarna i kryterium Tarskiego--Vaughta}
\begin{definicja}[Podstruktura elementarna]\label{def:model-elementary-substructure}
Niech $\mathfrak A\subseteq\mathfrak B$ będą strukturami
tego samego języka $\mathcal L$ o uniwersach $A\subseteq B$.
Podstrukturę $\mathfrak A$ nazywamy \emph{elementarną}
w $\mathfrak B$ i piszemy $\mathfrak A\preccurlyeq\mathfrak B$,
jeżeli dla każdego $n\ge0$, każdej formuły języka $\mathcal L$
o wszystkich zmiennych wolnych wśród parami różnych
$x_1,\ldots,x_n$ oraz wszystkich $a_1,\ldots,a_n\in A$
zachodzi
\[
\mathfrak A\models\varphi[a_1,\ldots,a_n]
\quad\Longleftrightarrow\quad
\mathfrak B\models\varphi[a_1,\ldots,a_n].
\]
Zapis $\mathfrak A\models\varphi[a_1,\ldots,a_n]$
oznacza spełnianie przy dowolnym wartościowaniu przypisującym
$x_i$ wartość $a_i$. Wynik nie zależy od wartości pozostałych
zmiennych, bo nie występują one wolno w $\varphi$
(Twierdzenie~\ref{thm:sem-coincidence}).
Dla $n=0$ chodzi o prawdziwość zdań.
\end{definicja}

\begin{uwaga}[Jak interpretować elementarność?]
Podstruktura zachowuje działania i relacje dotyczące swoich
elementów. Podstruktura \emph{elementarna} zachowuje ponadto
prawdziwość dowolnych formuł z parametrami z $A$.
Jeśli w większej strukturze istnieje element spełniający
formułę z tymi parametrami, w mniejszej strukturze musi
również istnieć odpowiedni świadek. Może to być inny element,
ale musi należeć do $A$.
\end{uwaga}

\begin{twierdzenie}[Kryterium Tarskiego--Vaughta]\label{thm:model-tarski-vaught}
Niech $\mathfrak A\subseteq\mathfrak B$ będą strukturami
tego samego języka $\mathcal L$ o uniwersach $A\subseteq B$.
Następujące warunki są równoważne:
\begin{enumerate}[label=(\roman*)]
\item $\mathfrak A\preccurlyeq\mathfrak B$;
\item dla każdej formuły $\psi(x,\bar y)$ języka $\mathcal L$,
której wszystkie zmienne wolne należą do listy
$x,\bar y$ parami różnych zmiennych, oraz dla każdej krotki
$\bar a$ elementów $A$ odpowiedniej długości,
jeśli istnieje $b\in B$ takie, że
$\mathfrak B\models\psi[b,\bar a]$,
to istnieje $a\in A$ takie, że
$\mathfrak B\models\psi[a,\bar a]$.
\end{enumerate}
\end{twierdzenie}
\begin{proof}
Załóżmy (i). Jeżeli $\mathfrak B\models\psi[b,\bar a]$
dla pewnego $b\in B$, to przy wartościach $\bar a\in A$
w $\mathfrak B$ spełniona jest formuła
$\exists x\,\psi(x,\bar y)$. Jej zmienne wolne należą do
listy $\bar y$, a zatem elementarność daje spełnianie
tej samej formuły w $\mathfrak A$ przy $\bar a$.
Istnieje więc $a\in A$ z
$\mathfrak A\models\psi[a,\bar a]$.
Ponowne użycie elementarności, teraz do $\psi$ i
parametrów $(a,\bar a)$, daje
$\mathfrak B\models\psi[a,\bar a]$.

Załóżmy (ii). Indukcyjnie po budowie formuły $\varphi$
udowodnimy, że dla każdego wartościowania $s$ w
$\mathfrak A$, traktowanego również jako wartościowanie
w $\mathfrak B$, zachodzi
$\mathfrak A,s\models\varphi\iff\mathfrak B,s\models\varphi$.
Wystarczy porównywać wartości zmiennych wolnych, ponieważ
spełnianie nie zależy od pozostałych wartości
(Twierdzenie~\ref{thm:sem-coincidence}).
Dla termów zgodność interpretacji wynika indukcyjnie z
definicji podstruktury, a dla atomów relacyjnych z
$R^{\mathfrak A}=R^{\mathfrak B}\cap A^k$,
gdzie $k$ jest arnością symbolu $R$.
Otrzymujemy zatem równoważność dla atomów równościowych
i relacyjnych. Negacja i spójniki dwuargumentowe
zachowują ją na mocy założenia indukcyjnego.

Niech $\varphi=\exists x\,\psi$. Wybieramy listę
$\bar y$ wszystkich pozostałych zmiennych wolnych
formuły $\psi$. Zmienna $x$ jest różna od zmiennych
z $\bar y$; w razie potrzeby dopuszczalnie przemianowujemy
zmienną związaną, korzystając z
Twierdzenia~\ref{thm:sem-alpha-renaming-xlv}.
Ustalmy $\bar a\in A^{|\bar y|}$.
Jeśli $\mathfrak A\models(\exists x\,\psi)[\bar a]$,
to istnieje $a\in A$ takie, że
$\mathfrak A\models\psi[a,\bar a]$.
Hipoteza indukcyjna daje
$\mathfrak B\models\psi[a,\bar a]$,
a zatem $\mathfrak B\models(\exists x\,\psi)[\bar a]$.
Odwrotnie, jeżeli $\mathfrak B\models(\exists x\,\psi)[\bar a]$,
to z (ii) otrzymujemy $a\in A$ spełniające
$\mathfrak B\models\psi[a,\bar a]$.
Hipoteza indukcyjna przenosi to na $\mathfrak A$,
więc $\mathfrak A\models(\exists x\,\psi)[\bar a]$.
Dla kwantyfikatora ogólnego korzystamy z semantycznej
równoważności
$\forall x\,\psi\iff\neg\exists x\,\neg\psi$:
przypadki negacji i kwantyfikatora egzystencjalnego
zostały już sprawdzone. Indukcja jest zakończona,
co daje (i).
\end{proof}

\begin{przyklad}[Dlaczego zwykła podstruktura nie musi być elementarna]
Wróćmy do $\mathfrak A=(\N,S)$ i $\mathfrak B=(\Z,S)$.
Zbiór $\N$ jest podstrukturą $\Z$, lecz nie jest
podstrukturą elementarną. Dla parametru $a=0\in\N$
w większej strukturze zachodzi $\exists x\,(S(x)=a)$,
ale nie ma takiego świadka w $\N$.
Właśnie tutaj zawodzi warunek Tarskiego--Vaughta.
\end{przyklad}

\section{Teoria struktury i równoważność elementarna}
\begin{definicja}[Równoważność elementarna]\label{def:model-elementary-equivalence}
Niech $\mathfrak A,\mathfrak B$ będą strukturami tego samego
języka pierwszego rzędu $\mathcal L$.
Dla dowolnej struktury $\mathfrak C$ języka $\mathcal L$ przez
$\operatorname{Th}(\mathfrak C)$ rozumiemy zbiór wszystkich zdań
języka $\mathcal L$ prawdziwych w $\mathfrak C$;
w Definicji~\ref{def:compactness-theory-n} rozważaliśmy
szczególny przypadek $\operatorname{Th}(\mathfrak N)$.
Mówimy, że $\mathfrak A$ i $\mathfrak B$ są
\emph{elementarnie równoważne}, i piszemy
$\mathfrak A\equiv\mathfrak B$, jeśli każde zdanie języka
$\mathcal L$ jest prawdziwe w $\mathfrak A$ dokładnie wtedy,
gdy jest prawdziwe w $\mathfrak B$, czyli
$\operatorname{Th}(\mathfrak A)=\operatorname{Th}(\mathfrak B)$.
\end{definicja}

\begin{twierdzenie}[Izomorfizm i równoważność elementarna]\label{thm:model-iso-elementary}
Niech $\mathfrak A$ i $\mathfrak B$ będą strukturami tego
samego języka pierwszego rzędu $\mathcal L$.
Jeżeli $\mathfrak A\cong\mathfrak B$, to
$\mathfrak A\equiv\mathfrak B$.
Istnieją jednak elementarnie równoważne struktury
pierwszego rzędu, które nie są izomorficzne.
\end{twierdzenie}
\begin{proof}
Pierwsze stwierdzenie wynika z
Twierdzenia~\ref{thm:model-iso-invariance} zastosowanego do zdań.
Dla drugiego wykorzystajmy Twierdzenie~\ref{thm:compactness-nonstandard}.
Niech $\mathfrak M^+$ oznacza otrzymany tam model języka
$\mathcal L_{\mathrm{ar}}\cup\{c\}$, a
$\mathfrak M=\mathfrak M^+\upharpoonright\mathcal L_{\mathrm{ar}}$
jego redukt. Z $\mathfrak M\models\operatorname{Th}(\mathfrak N)$
wynika $\mathfrak M\equiv\mathfrak N$: jeśli
$\mathfrak N\not\models\sigma$, to
$\mathfrak N\models\neg\sigma$, a więc i
$\mathfrak M\models\neg\sigma$.

Gdyby istniał izomorfizm $h:\mathfrak M\to\mathfrak N$,
obraz elementu $a:=c^{\mathfrak M^+}$ byłby równy
$m=h(a)\in\N$. Chociaż $c$ nie jest symbolem języka reduktu,
$a$ pozostaje elementem jego uniwersum.
Z Twierdzenia~\ref{thm:compactness-nonstandard} mamy
$\overline m^{\mathfrak M}<^{\mathfrak M}a$.
Izomorfizm zachowuje $0$, $S$ i $<$, toteż przenosi tę
nierówność na $m<m$ w $\mathfrak N$, co jest sprzecznością.
\end{proof}

\begin{uwaga}[Trzy poziomy podobieństwa]
Relacja $\mathfrak A\subseteq\mathfrak B$ wymaga zgodności
interpretacji symboli na $A$.
Warunek $\mathfrak A\preccurlyeq\mathfrak B$ zapewnia ponadto
zachowanie wszystkich formuł z parametrami z $A$.
Natomiast $\mathfrak A\equiv\mathfrak B$ oznacza zgodność
prawdziwości zdań i nie wymaga ani zawierania uniwersów,
ani istnienia odwzorowania między strukturami.
\end{uwaga}